\section{Prepared Statement dengan MySQLi: Mencegah SQL Injection}

\textbf{SQL injection} adalah serangan di mana penyerang menyisipkan kode SQL berbahaya melalui input pengguna. Misalnya, input \code{' OR 1=1 --} pada form login dapat melewati autentikasi jika query dibangun dengan penggabungan string. \textbf{Prepared statement} mencegah serangan ini dengan memisahkan template query dari data: MySQL memperlakukan data sebagai literal (nilai), bukan sebagai bagian dari perintah SQL \cite{phpRightWay}.

\begin{contoh}
\textbf{Studi Kasus: Pencegahan SQL Injection}

\textbf{Query rentan}: \code{"SELECT * FROM users WHERE user='\$u' AND pass='\$p'"}. Jika \code{\$u = "admin' --"}, query menjadi: \code{SELECT * FROM users WHERE user='admin' --' AND pass='...'} --- komentar SQL (\code{--}) menghilangkan cek password.

\textbf{Dengan prepared statement}: \code{\$stmt = \$conn->prepare("SELECT * FROM users WHERE user=? AND pass=?")}. Input \code{"admin' --"} diperlakukan sebagai string literal, bukan perintah SQL. Serangan gagal.
\end{contoh}

\begin{figure}[H]
\centering
\begin{tikzpicture}[
  scale=0.88,
  transform shape,
  box/.style={draw, rounded corners, minimum width=2.75cm, minimum height=0.9cm, align=center, font=\footnotesize, fill=#1},
  arrow/.style={-Stealth, thick},
  node distance=0.85cm
]
  \node[box=blue!15] (prep) {1. Prepare\\template query};
  \node[box=green!15, right=1.15cm of prep] (bind) {2. Bind\\parameter};
  \node[box=orange!20, right=1.15cm of bind] (exec) {3. Execute};
  \node[box=purple!15, right=1.15cm of exec] (fetch) {4. Fetch\\hasil};

  \draw[arrow] (prep) -- (bind);
  \draw[arrow] (bind) -- (exec);
  \draw[arrow] (exec) -- (fetch);
\end{tikzpicture}
\caption{Alur prepared statement: prepare, bind, execute, fetch}
\end{figure}

\begin{webcode}[language=PHP, caption={Prepared statement MySQLi lengkap}]
<?php
// Langkah 1: Prepare - template query dengan placeholder ?
$stmt = $conn->prepare(
  "SELECT id, nama, email FROM pengguna WHERE email = ? AND aktif = ?"
);

// Langkah 2: Bind - tipe: s=string, i=integer, d=double, b=blob
$email = $_POST['email'];
$aktif = 1;
$stmt->bind_param("si", $email, $aktif);

// Langkah 3: Execute
$stmt->execute();

// Langkah 4: Fetch hasil
$result = $stmt->get_result();
while ($row = $result->fetch_assoc()) {
  echo htmlspecialchars($row['nama']) . "\n";
}

$stmt->close();
?>
\end{webcode}

